\section{Systemer af lineære ODE'er}
En ordinær differentialligning er en ligning på formen
\[
F(f^{(n)}(t), \cdots, f'(t), f(t), t) = 0 \kildetag{Definition 13.0.1}
\]
Notationen betyder at $F$ er en funktion der tager $n+2$ variabler, og de variabler er så de forskellige afledede af funktionen $f(t)$ samt $t$ selv.

En løsning er en funktion $f(t)$, som opfylder ligningen for alle $t \in \mathbb{R}$.

\subsection{Lineær førsteordens ODE}

\[
f'(t) = a(t) f(t) + q(t)
\]
hvor $a(t), q(t)$ er funktioner i $t$ og $q(t)$ er en \emph{forcing function} er den generelle form for en lineær førsteordens ODE.

Hvis $q(t) = 0$, er ODE'en homogen.

\subsection{Er differentialligningen lineær?}
\label{sec:diff.lineær}
En differentialligning med $f(t)$ er lineær hvis den kan skrives på formen:
\[
f'(t) + a(t) \cdot f(t) = q(t)
\]

Hver af $f^{(n)}(t)$ skal have eksponent $1$ og må ikke ganges med andre $f^{(n)}(t)$.

Ingen af $f^{(n)}(t)$ må være i en funktion som $\sin(\cdot)$, $e^{(\cdot)}$ osv.

\subsection{Find førstegradspolynomium}
Givet en lineær førsteordens ODE og viden at løsningen er et polynomium af første grad, $f(t) = c_1 t + c_0$, kan man finde konstanterne $c_1$ og $c_0$ ved at indsætte $f(t)$ i ODE'en og løse for konstanterne.

$f'(t) = c_1$, så indsætter vi i ODE'en og får en ligning i $t$. Vi kan så sammenligne koefficienterne for hver potens af $t$ på begge sider af ligningen for at finde $c_1$ og $c_0$.

\subsection{Løs matrix ODE}
Givet en matrix $A$ og en ODE på formen
\[
\begin{bmatrix}
    f_1'(t) \\
    f_2'(t) \\
\end{bmatrix}
=
\textbf{A} \cdot
\begin{bmatrix}
    f_1(t) \\
    f_2(t) \\
\end{bmatrix}
\]
Kan løsningen findes ved at finde egenværdierne og egenvektorerne for matrixen $\mathbf{A}$.
Da kan løsningen skrives som
\[
\begin{split}
    f_1(t) \\
    f_2(t) \\
\end{split}
= c_1 e^{\lambda_1 t} \vec{v_1} + c_2 e^{\lambda_2 t} \vec{v_2}
\]

\subsection{Løs for begyndelsesværdier}
Givet begyndelsesværdier og den generelle løsning til en ODE
\[
\begin{bmatrix}
    f_1(t) \\
    f_2(t) \\
\end{bmatrix}
=
c_1 \cdot
\begin{bmatrix}
    y_1 \\
    y_2 \\
\end{bmatrix}
\cdot e^{\lambda_1 t}
+ c_2 \cdot
\begin{bmatrix}
    z_1 \\
    z_2 \\
\end{bmatrix}
\cdot e^{\lambda_2 t}
\]

Kan $c_1, c_2$ findes ved at udregne LLS'et med indsatte begyndelsesværdier.

\eksempel{
Egenværdier $2, 4$ med egenvektorer 
$\begin{bmatrix}
    1 \\
    1
\end{bmatrix}$,
$\begin{bmatrix}
    1 \\
    -1
\end{bmatrix}$

Den generelle løsning:
\[
\begin{bmatrix}
    f_1(t) \\
    f_2(t) \\
\end{bmatrix}
=
c_1 \cdot
\begin{bmatrix}
    1 \\
    1 \\
\end{bmatrix}
\cdot e^{2 t}
+ c_2 \cdot
\begin{bmatrix}
    1 \\
    -1 \\
\end{bmatrix}
\cdot e^{4 t}
\]
Med begyndelsesværdier
\(
\begin{bmatrix}
    f_1(0) \\
    f_2(0)
\end{bmatrix}
= \begin{bmatrix}
    6 \\
    0
\end{bmatrix}
\)
Vil værdierne indsættes og give
\[
\begin{bmatrix}
    6 \\
    0 \\
\end{bmatrix}
=
\begin{bmatrix}
    c_1 \cdot e^{0} + c_2 \cdot e^{0} \\
    c_1 \cdot e^{0} - c_2 \cdot e^{0} \\
\end{bmatrix}
\xrightarrow{\text{Siden } e^0=1}
\begin{bmatrix}
    6 \\
    0 \\
\end{bmatrix}
=
\begin{bmatrix}
    c_1 + c_2 \\
    c_1 - c_2 \\
\end{bmatrix}
\]
Som LLS er det:
\[
\begin{cases}
6 = c_1 + c_2 \\
0 = c_1 - c_2
\end{cases}
\]

Her er det tydeligt at $c_1 = c_2$, og at $c_1 + c_2 = 6$, derfor $c_1 = c_2 = 3$.
\[
\begin{bmatrix}
    f_1(t) \\
    f_2(t) \\
\end{bmatrix}
=
3 \cdot
\begin{bmatrix}
    1 \\
    1 \\
\end{bmatrix}
\cdot e^{2 t}
+ 3 \cdot
\begin{bmatrix}
    1 \\
    -1 \\
\end{bmatrix}
\cdot e^{4 t}
\]
}


\newpage

\subsection{Løs Lineære anden-ordens ODE'er}
\subsubsection{Homogen}
En homogen lineær anden-ordens ODE er på formen
\[
f''(t) + a_1(t) f'(t) + a_0(t) f(t) = 0
\]

\textbf{Karakteristisk Polynomium}

Det karakteristiske polynomium for en homogen lineær anden-ordens ODE med konstante koefficienter er:
\[
Z^2 + a_1 Z + a_0 = 0
\]

Find determinanten af det karaktistiske polynomium:
\[
D = {a_1}^2 - 4 a_0
\]

\hrule
\

\textbf{D > 0}

Her vil rødderne være:
\[
\lambda_1 = \frac{-a_1 + \sqrt{D}}{2} \qquad \lambda_2 = \frac{-a_1 - \sqrt{D}}{2}
\] 

Den generelle løsning er da:
\[
f(t) = c_1 e^{\lambda_1 t} + c_2 e^{\lambda_2 t}
\]

\hrule
\

\textbf{D = 0}

Her vil rødderne være:
\[
\lambda_1 = \lambda_2 = \frac{-a_1}{2}
\]

Den generelle løsning er da:
\[
c_1 e^{\lambda_1 t} + c_2 t e^{\lambda_2 t} 
\]

\hrule
\

\textbf{D < 0} 

Her vil rødderne være:
\[
\lambda_1 = \frac{-a_1 + i \sqrt{|D|}}{2} \qquad \lambda_2 = \frac{-a_1 - i \sqrt{|D|}}{2}
\]

Den generelle (reelle) løsning er da:
\[
f(t) = c_1 e^{\frac{-a_1}{2} t} \cos\left(\frac{\sqrt{|D|}}{2} t\right) + c_2 e^{\frac{-a_1}{2} t} \sin\left(\frac{\sqrt{|D|}}{2} t\right)
\]
Eller
\[
f(t) = e^{\frac{-a_1}{2} t} \left( c_1 \cos\left(\frac{\sqrt{|D|}}{2} t\right) + c_2 \sin\left(\frac{\sqrt{|D|}}{2} t\right) \right)
\]

\subsection{Lineartitet}

En ODE som:
\[
\begin{split}
    f_1'(t) = f_1(t) + 2 f_2(t) \\
    f_2'(t) = 2 f_1(t) + f_2(t) \\
\end{split}
\]
Er lineær, da hver af ligningerne er lineær som beskrevet i sektion \ref{sec:diff.lineær}.